\section{The reflection extension and its outer actions}
\label{ext:section}

The models $M=M_0$ and $N=M_\ell$ realize two graph subgroups of
$\mathsf A\times\Ftwo$ on $Q_2$.  We now place these actions inside
$\Out(\Pi)$, where $\Pi=\pi_1(Q_2)$.  A reflection extension of
$\Pi$ will identify the prohibited cyclic action and will also
distinguish the two graph subgroups by an intrinsic property of their
elements.  Throughout this section we identify fundamental groups
with deck groups using left actions.

\subsection{The lifted development}

Let $\widetilde P\to P$ be the universal cover of the chamber.  Its
deck group is $H$.  The preimage of $S$ is connected because
$\pi_1(S)\to H$ is surjective; denote it by $\widehat S$.  It is the
regular cover of $S$ associated to the kernel of that epimorphism,
and its deck group is $H$.  Let $\widehat{\mathcal L}$ be the lift of
the chosen PL triangulation $\mathcal L$, and write
$\widehat{\mathcal V}$ for its vertex set.  For
$u\in\widehat{\mathcal V}$, write $\overline u\in\mathcal V$ for its
image.

\begin{lemma}\label{ext:lifted-nerve}
The complex $\widehat{\mathcal L}$ is a connected, locally finite,
flag-no-square PL triangulation of the three-manifold $\widehat S$.
The deck action of $H$ is free on its vertices, and its quotient is
$\mathcal L$.  In particular, each vertex link is a PL two-sphere.
\end{lemma}

\begin{proof}
The local assertions about the triangulation and its links follow
by lifting PL charts.  The action is free because a deck
transformation of a connected covering that fixes a point is the
identity.  Its quotient is $\mathcal L$ because the covering is
regular with deck group $H$.

A clique upstairs projects to a clique of distinct vertices
downstairs.  Flagness gives a simplex on their images.  Lift this
simplex from one vertex of the clique.  Uniqueness of the lifts of
its incident edges places every other vertex of the clique in that
lifted simplex.  Thus $\widehat{\mathcal L}$ is flag.

For a four-cycle upstairs, adjacent vertices have distinct images.
Opposite vertices also have distinct images: otherwise, at an
intervening vertex, two incident lifted edges would be lifts of the
same downstairs edge.  The projected four-cycle therefore has four
distinct vertices.  It has a diagonal because $\mathcal L$ has no
induced four-cycle.  The diagonal and two consecutive sides span a
triangle downstairs by flagness.  Lifting that triangle from their
common vertex lifts the diagonal to the original four-cycle.
Consequently the lifted cycle is not induced.
\end{proof}

Let $W$ be the right-angled Coxeter group on the graph of
$\widehat{\mathcal L}$:
\[
 W=\left\langle u\ (u\in\widehat{\mathcal V})\ \middle|\
 u^2=1,\quad uv=vu\ \text{if }uv\text{ is an edge}\right\rangle .
\]
We use a vertex and its Coxeter generator interchangeably.  The
deck action permutes these generators and gives the semidirect
product
\[
 \mathcal E=W\rtimes H,\qquad
 (d,h)(d',h')=(d\,h(d'),hh').
\]
Define an $H$-invariant homomorphism
$\chi_0\colon W\to\mathsf A_0$ by $\chi_0(u)=e_{\overline u}$.
Extend it to $\mathcal E$ by making it zero on $H$, and let
$\chi\colon\mathcal E\to\mathsf A$ be its composition with the
quotient $\mathsf A_0\to\mathsf A$.  Finally, define
\begin{equation}\label{ext:eta}
 \eta\colon\mathcal E\longrightarrow\mathsf A\times\Ftwo,\qquad
 \eta(d,h)=\bigl(\chi(d,h),w(h)\bmod 2\bigr).
\end{equation}
This map is onto: the generators of $W$ give every label $a_s$,
and the epimorphism $w\colon H\to\mathbb Z$ gives the second factor
independently.

The stabilizers in the development will be special subgroups.
We record their elementary support properties before using them.

\begin{lemma}[Special subgroups and support]\label{ext:support}
Let $W(\mathcal G)$ be a right-angled Coxeter group on an arbitrary
simple graph with vertex set $V$.  For $J\subseteq V$, the special
subgroup $W_J$ generated by $J$ is the right-angled Coxeter group on the induced
graph, and
\[
 W_J\cap W_K=W_{J\cap K}.
\]
Every $d\in W(\mathcal G)$ has a unique smallest vertex set
$\supp(d)$ for which $d\in W_{\supp(d)}$.  This support is
finite, and graph automorphisms transport supports.
\end{lemma}

\begin{proof}
Killing every generator outside $J$ defines a retraction onto
the Coxeter group of the induced graph on $J$.  Its composition
with the natural map of that group into $W(\mathcal G)$ is the
identity, so that natural map is injective and gives $W_J$.
If $d\in W_J\cap W_K$, apply the retraction onto $W_J$ to a word
for $d$ on $K$.  The result is a word on $J\cap K$ for $d$,
proving the intersection formula.

Choose a finite set $F$ supporting a word for $d$.  Intersect
$F$ with every supporting set.  Because $F$ is finite, the
result is already the intersection of finitely many of those
sets; repeated use of the intersection formula shows that it
still supports $d$.  It is therefore the unique smallest
support.  Its characterization makes it equivariant under graph
automorphisms.  This is the right-angled form of the standard
Coxeter support statement
\cite[Proposition~4.1.1 and Corollary~4.1.2]{Davis2008}.
\end{proof}

Apply this notation to $W$.  At $x\in\widetilde P$, let
$J(x)$ be the vertices of the lifted panels containing $x$, with
$J(x)=\varnothing$ off $\widehat S$.  This set spans a simplex, so
$W_{J(x)}$ is the elementary abelian group on $J(x)$.  Form
\begin{equation}\label{ext:universal-space}
 \widetilde Q=(W\times\widetilde P)/\sim,\qquad
 (c,x)\sim(c',x)
 \ \Longleftrightarrow\ c^{-1}c'\in W_{J(x)}.
\end{equation}
The relation never changes the chamber coordinate.  Since
$J(hx)=hJ(x)$, the formula
\begin{equation}\label{ext:action}
 (d,h)[c,x]=[d\,h(c),hx]
\end{equation}
defines an action of $\mathcal E$ on $\widetilde Q$.

\begin{proposition}\label{ext:development}
The space $\widetilde Q$ is contractible.  If $\overline x$ and
$x_2$ are the images of $x$ in $P$ and $P_2$, respectively, the
maps
\[
 \begin{aligned}
 [c,x]&\longmapsto[\chi_0(c),\overline x]\in Q^0,\\
 [c,x]&\longmapsto[\chi(c),\overline x]\in Q,\\
 [c,x]&\longmapsto[\chi(c),x_2]\in Q_2
 \end{aligned}
\]
are universal coverings.  With the chosen deck-group convention,
\begin{equation}\label{ext:groups}
 \pi_1(Q^0)=\ker\chi_0,\qquad
 G:=\pi_1(Q)=\ker\chi,\qquad
 \Pi:=\pi_1(Q_2)=\ker\eta .
\end{equation}
The homomorphism $w\colon G\to\mathbb Z$ induced by the folding map
is $w$ composed with the projection $\mathcal E\to H$.
\end{proposition}

\begin{proof}
The contractibility assertion is the universal-development
argument in \cite[Section~3, subsection \emph{The universal
development}, immediately before Proposition~3.2]{PD}.  Its
chamber hypotheses are the ones verified for $(P,S)$ in the
reflection construction.  That argument filters the chambers
of \eqref{ext:universal-space} by word length in $W$.  A new
chamber is contractible and meets the preceding chambers in
the nonempty ball of its descent panels.  The argument treats an
infinite word-length layer simultaneously and takes a CW direct
limit, giving a contractible development even when the lifted
generator set is infinite.  This is the lifted-mirror construction
of Davis \cite[Section~9.1, Equations~(9.1)--(9.3) and
Theorem~9.1.5]{Davis2008}; Equation~(9.3) there is
\eqref{ext:action}.

Here is the quotient and covering check for the label groups in
this paper.  The set $J(x)$ projects bijectively to the set of
panels meeting $\overline x$.  The restrictions of $\chi_0$ and
$\chi$ to $W_{J(x)}$ are therefore injective: for $\chi_0$ the
images are distinct coordinate vectors, and for $\chi$ they are
the simplex-independent labels $a_s$.  Their images are exactly
the label subgroups used in the respective reflection
identifications.  This proves that the three displayed maps are
well-defined.  They are onto because every label has a
representative in $W$ and every chamber coordinate lifts.

We check their fibers uniformly.  Use $\lambda=\chi_0$ or $\chi$
and the chamber cover $P_d$, where the three cases are
$(\lambda,d)=(\chi_0,1),(\chi,1),(\chi,2)$, with
$P_1=P$, $H_1=H$, and $H_2=w^{-1}(2\mathbb Z)$.
For these cases put
\[
 K_{\lambda,d}
   =\{(d_0,h)\in\mathcal E:\lambda(d_0)=0,\ h\in H_d\}.
\]
Suppose $[c,x]$ and $[c',x']$ have the same image.  Equality of
their chamber coordinates gives $x'=hx$ for some $h\in H_d$,
and equality of the reflected points says that
$\lambda(c')-\lambda(c)$ belongs to
$\lambda(W_{J(x')})$.  Choose $t\in W_{J(x')}$ with
\[
 \lambda(t)=\lambda(c')-\lambda(c).
\]
Then
$d_0=c't\,h(c)^{-1}$ satisfies $\lambda(d_0)=0$, since
$\lambda$ is $H$-invariant and the label groups have
characteristic two.  Formula \eqref{ext:action} sends $[c,x]$
by $(d_0,h)$ to $[c't,x']=[c',x']$.  Conversely, every element
of $K_{\lambda,d}$ preserves the displayed map.  Thus its
fibers are exactly the $K_{\lambda,d}$-orbits.

These orbit maps are covering maps, including at the panels.
At an interior point choose an evenly covered chamber
neighborhood disjoint from the boundary.  At a boundary point
choose a small lifted collar neighborhood meeting only the panels
in $J(x)$ and mapping homeomorphically to its chamber
neighborhood.  The adjacent chamber sectors are indexed by
$W_{J(x)}$ upstairs and by $\lambda(W_{J(x)})$ downstairs.
The isomorphism between these two finite groups identifies their
sector gluings, so each such assembled neighborhood maps
homeomorphically onto the corresponding neighborhood in the
reflection.  The different $K_{\lambda,d}$-translates give the disjoint
sheets.  The point stabilizer in $\mathcal E$ at
$[c,x]$ is $cW_{J(x)}c^{-1}$: a fixing element has $h=1$ because
$H$ acts freely on $\widetilde P$, and the relation
\eqref{ext:universal-space} then gives this conjugate.  The subgroup
$K_{\lambda,d}$ meets it trivially by the injectivity just proved.
These are precisely the covering charts used in the cited
universal-development argument.

For comparison with the notation of \cite{PD}, let
$C_{\mathcal L}$ be the downstairs right-angled Coxeter group.
The homomorphism forgetting the lift of each generator is
$q\colon W\to C_{\mathcal L}$, and that paper uses
$\Delta=\ker(C_{\mathcal L}\to\mathsf A_0)$ and
$\rho(d,h)=q(d)$.  Its group $\rho^{-1}(\Delta)$ is exactly
$\ker\chi_0$.  Thus the first identification in
\eqref{ext:groups} is also the group formula in the cited
subsection of \cite{PD}.  The fiber calculation gives the
identifications for $Q$ and $Q_2$.  Since $\widetilde Q$ is
contractible, these coverings are universal.

Finally, the map $[c,x]\mapsto x$ from $\widetilde Q$ to
$\widetilde P$ is equivariant for the projection
$\mathcal E\to H$.  It covers the folding map $Q\to P$.
Consequently the fold-induced homomorphism on $G$ is that
projection, which proves the assertion about $w$.
\end{proof}

The preimages of the two coordinate factors are
\begin{equation}\label{ext:e2}
 \begin{aligned}
 G&=\eta^{-1}(\{0\}\times\Ftwo),\\
 \mathcal E_2&=W\rtimes H_2
   =\eta^{-1}(\mathsf A\times\{0\}),\\
 G\cap\mathcal E_2&=\Pi.
 \end{aligned}
\end{equation}
All three groups are normal in $\mathcal E$.

\subsection{Embedding the extension in the automorphism group}

The extension acts on its normal subgroup $\Pi$ by conjugation.
Finite support and the free deck action show that this action
is faithful.

\begin{proposition}\label{ext:embedding}
Conjugation gives an injection
$\mathcal E\hookrightarrow\Aut(\Pi)$.  It identifies $\Pi$ with
$\Inn(\Pi)$ and induces an inclusion
\begin{equation}\label{ext:outer}
 \mathcal E/\Pi
   =\mathsf A\times\langle\delta\rangle
   \ \leq\ \Out(\Pi),
 \qquad \delta=(0,1).
\end{equation}
Here $\langle\delta\rangle\cong\Ftwo$.  Inside $\Aut(\Pi)$,
$\mathcal E$ is the full preimage of the product in
\eqref{ext:outer}; $\mathcal E_2$ and $G$ are the full preimages of
its factors $\mathsf A$ and $\langle\delta\rangle$, respectively.
\end{proposition}

\begin{proof}
The subgroup $\Pi$ is normal in $\mathcal E$ by
\eqref{ext:groups}.  Suppose that $(d,h)$ centralizes $\Pi$.
It commutes with $(1,k)$ for every $k\in H_2\subseteq\Pi$.
The semidirect-product formula gives
\[
 hk=kh,\qquad d=k(d)\qquad(k\in H_2).
\]
The finite-index subgroup $H_2$ of the nonabelian free group $H$
is itself nonabelian.  If $h\ne1$, the first equality would put
$H_2$ inside the cyclic centralizer of $h$ in $H$, a
contradiction.  Thus $h=1$.

The second equality and Lemma~\ref{ext:support} make the finite
set $\supp(d)$ invariant under $H_2$.  Every $H_2$-orbit of
vertices is infinite: $H_2$ is infinite and its deck action is
free.  Hence $\supp(d)$ is empty, so $d=1$.  This proves the
injection and also proves that $Z(\Pi)=1$.

Conjugation by the elements of $\Pi$ gives exactly
$\Inn(\Pi)$, and it is injective because the center is trivial.
Taking quotients now gives \eqref{ext:outer} through
\eqref{ext:eta}.  To see that the preimage is full, let an
automorphism of $\Pi$ have the same outer class as conjugation
by some $e\in\mathcal E$.  It differs from that conjugation by
conjugation by an element $\gamma\in\Pi$, and is therefore
conjugation by $\gamma e\in\mathcal E$.  Taking the preimages of
the two factors and using \eqref{ext:e2} gives the last
assertion.
\end{proof}

\begin{proposition}\label{ext:outer-actions}
Under \eqref{ext:outer}, the action of $\mathcal E/\Pi$ on $Q_2$
is the label-translation and cyclic-deck action
\[
 (a,t)[b,y]=[a+b,\theta^t y],
 \qquad (a,t)\in\mathsf A\times\Ftwo.
\]
With the markings supplied by Proposition~\ref{mod:models}, the
graph action on $M_\epsilon$ from Proposition~\ref{mod:actions}
induces exactly the subgroup
\[
 \mathsf B_\epsilon
   =\{(a,\epsilon(a)):a\in\mathsf A\}
   \leq\Out(\Pi).
\]
In particular, every element of $\mathsf A=\mathsf B_0$ is
induced by a homeomorphism of $M$, and every element of
$\mathsf B_\ell$ is induced by a homeomorphism of $N$.
\end{proposition}

\begin{proof}
In the covering map to $Q_2$ from
Proposition~\ref{ext:development}, Equation~\eqref{ext:action}
changes the label by $\chi(d,h)$ and the chamber coordinate by
the action of $h$ on $P_2$.  The latter action is
$\theta^{w(h)\bmod 2}$.  This gives the displayed formula.
Conjugation on the universal deck group is the induced outer
automorphism, so \eqref{ext:outer} identifies these actions with
the stated outer classes.  The equivariance of the model maps
in Proposition~\ref{mod:actions} then identifies the classes of
the graph actions on the manifolds.
\end{proof}

\begin{proposition}\label{ext:delta}
With the marking $\pi_1(M)\cong\Pi$, no homeomorphism of $M$
induces the outer automorphism $\delta$.
\end{proposition}

\begin{proof}[Proof of Proposition~\ref{ext:delta} and
Theorem~\ref{thm:self-map}]
The quotient reflection $Q$ is a connected finite aspherical
complex.  Its folding homomorphism is the epimorphism
$w\colon G\to\mathbb Z$ of Proposition~\ref{ext:development}.
The space $Q_2$ has the closed topological four-manifold model
$M$ by Proposition~\ref{mod:models}, whereas
Proposition~\ref{mod:odd-models} excludes such a model for every
odd $d\geq3$.  These facts verify all hypotheses of
Theorem~\ref{cyc:criterion}, with $n=4$.

Choose $h\in H$ with $w(h)=1$.  The element
$\gamma=(1,h)$ lies in $G=\ker\chi$, has $w(\gamma)=1$, and
maps to $(0,1)$ under \eqref{ext:eta}.  Consequently conjugation
by $\gamma$ on $\Pi$ represents $\delta$ in
\eqref{ext:outer}.  The cyclic-cover criterion gives a
self-homotopy equivalence $f\colon M\to M$ inducing this class,
and asserts that no homeomorphism of $M$ induces it.  This proves
Proposition~\ref{ext:delta}.  A map homotopic to $f$ induces the
same outer automorphism, so $f$ is not homotopic to a
homeomorphism.  This proves Theorem~\ref{thm:self-map}.
\end{proof}

\begin{proposition}\label{ext:nonconjugate}
The subgroups $\mathsf A$ and $\mathsf B_\ell$ of $\Out(\Pi)$
are not conjugate.
\end{proposition}

\begin{proof}
Consider the property that an outer automorphism has a
finite-order lift to $\Aut(\Pi)$.  This property is invariant
under conjugation in $\Out(\Pi)$: if $F$ represents the
conjugating outer automorphism and $\alpha$ is a finite-order
lift, then $F\alpha F^{-1}$ is a finite-order lift of the
conjugate.

Every lift of an element $(a,1)$ in \eqref{ext:outer} belongs
to $\mathcal E$ by Proposition~\ref{ext:embedding}.  Its
projection to $H$ has odd, hence nonzero, $w$-value.  It
therefore has infinite order in $\mathcal E$ and, by the
injection into $\Aut(\Pi)$, infinite order as an automorphism.
Thus no $(a,1)$ has the property.  In contrast, each
$(a_s,0)$ has a finite-order lift: any Coxeter generator above
$s$ is an involution in $\mathcal E$ lifting it.

The graph subgroup $\mathsf B_\ell$ and the horizontal subgroup
$\mathsf A$ have identical elements over $\ker\ell$.  Every
element of $\mathsf B_\ell$ over its complement has second
coordinate $1$ and lacks a finite-order lift.  Some $a_s$
lies in that complement, since the $a_s$ span $\mathsf A$ and
$\ell\ne0$; the corresponding element of $\mathsf A$ does have
such a lift.  Hence $\mathsf B_\ell$ has strictly fewer
elements with the property than $\mathsf A$.  Conjugate finite
subgroups have the same number, proving the claim.
\end{proof}

The normalizer calculation below will ensure that a homeomorphism
between the models would force the conjugacy just excluded.  It
will be combined with the finiteness of $\Out(\Pi)$ established in the next section.
